\documentclass[11pt]{article}
\usepackage[spanish]{babel}
\usepackage{amsmath}
\usepackage{amssymb}
\usepackage{graphicx}
\graphicspath{ {./images/}, {~/Apunte/images/} }
\usepackage[utf8]{inputenc}
\usepackage[T1]{fontenc}
\usepackage[margin=1in]{geometry}
\usepackage{fancyhdr}
\usepackage{hyperref}

\hypersetup{
    pdftitle={Trabajo Previo Laboratorio Python},
    pdfauthor={Felipe Colli Olea},
    colorlinks=true,
    linkcolor=black,
    urlcolor=cyan
}

\pagestyle{fancy}
\fancyhf{}
\fancyfoot[C]{\footnotesize Felipe Colli, Todos los Derechos Reservados}
\fancyfoot[R]{\thepage}
\renewcommand{\headrulewidth}{0pt}

\title{Preparación Laboratorio Python: \\ \large{IDLE, Funciones, Condicionales e Input}}
\author{Felipe Colli Olea}
\date{\today}

\begin{document}

\maketitle
\thispagestyle{fancy}

% ------------------------------------------------------------------------------
% FUNCIONES
% ------------------------------------------------------------------------------
\section{Funciones}

\textbf{Pregunta:} Al ejecutar el script, la última instrucción (\texttt{x = area(2,3)}) da un error. ¿Cuál es la razón? \\

\textbf{Respuesta:} El error de ejecución (\texttt{TypeError}) se produce debido a una discrepancia en el paso de argumentos. La función \texttt{area(x)} fue definida formalmente para recibir \textbf{exactamente un argumento posicional} (el valor \texttt{x}). Sin embargo, en la instrucción \texttt{area(2,3)}, se está intentando invocar la función enviándole \textbf{dos argumentos} independientes (\texttt{2} y \texttt{3}). El intérprete de Python detecta esta incongruencia en la firma de la función y detiene la ejecución.

\textit{Nota:} Cabe destacar que el script original contiene además un error de alcance de variables (\textit{scope}), ya que la función recibe el parámetro \texttt{x} pero intenta calcular \texttt{math.pi*r**2}, lo que generaría un \texttt{NameError} previo al intentar evaluar \texttt{r}.

% ------------------------------------------------------------------------------
% INGRESO DE VALORES
% ------------------------------------------------------------------------------
\section{Ingreso de valores por teclado}

\textbf{Pregunta:} ¿Cuál es el efecto que tiene la instrucción de ingreso (\texttt{input}) en el ejemplo completo? \\

\textbf{Respuesta:} La función \texttt{input()} tiene como efecto detener el hilo de ejecución del programa para esperar una interacción del usuario a través de la entrada estándar (\textit{standard input}). El dato capturado siempre es interpretado por defecto como una cadena de texto (\texttt{string}). 

Al anidar esta instrucción dentro de la función \texttt{int()} —formando \texttt{int(input(...))}—, el efecto logrado es un \textit{casting} (conversión de tipo de dato). Se transforma la cadena de texto a un número entero, lo cual es estrictamente necesario para que el intérprete pueda realizar las operaciones aritméticas posteriores del script (como \texttt{edad + 5}) sin generar un error de incompatibilidad de tipos (\texttt{TypeError}).

% ------------------------------------------------------------------------------
% EJECUCIÓN CONDICIONAL
% ------------------------------------------------------------------------------
\section{Ejecución condicional de instrucciones}

\textbf{Pregunta:} Al ingresar los valores 6 y 6. Comente por qué la respuesta entregada es ``el segundo es mayor que el primero''. \\

\textbf{Respuesta:} El resultado se debe a un fallo de diseño en la lógica condicional del algoritmo. Al ingresar \texttt{x = 6} e \texttt{y = 6}, la evaluación de la compuerta lógica del bloque \texttt{if (x > y)} se traduce como \texttt{6 > 6}. Matemáticamente esta proposición es \textbf{Falsa}.

Dado que la condición inicial no se cumple, y al carecer el programa de un bloque de control intermedio (como un \texttt{elif x == y:}) diseñado para manejar casos de igualdad, el flujo de ejecución decae automáticamente en la instrucción por defecto del bloque \texttt{else}. Por lo tanto, el intérprete asume ciegamente la instrucción contenida en dicho bloque, imprimiendo el texto ``el segundo es mayor que el primero'', a pesar de ser semánticamente incorrecto para el caso ingresado.

\end{document}
